\documentclass[12pt]{article}

\include{../style}

\begin{document}

\title{Cosc 30: Discrete Mathematics}

\author{Prishita Dharampal}
\date{}
\maketitle
\vspace{-0.3in}



\begin{problem}{1}
    \textbf{Better of the two.} 
    
    Toss a biased coin repeatedly until getting the first head, where the probability of landing on head is $p$ and on tail is $1-p$. Let $X$ be the number of tosses (including the one that lands on head). Now repeat the process independently a second time; denote by $Y$ the number of tosses to get the first head.

    \begin{enumerate}
        \item[(a)] What is the probability that $X = Y$? Justify your answer. (Your solution should contain $p$ as a parameter.)

        \item[(b)] Let $Z$ be any random variable defined on a sample space $S$, taking values in the nonnegative integers (that is, $Z : S \to \mathbb{N}_{\ge 0}$). Prove that
        \[
        E[Z] = \sum_{k \ge 1} \Pr[Z \ge k].
        \]

        \item[(c)] What is the expected value of $\min(X,Y)$? Justify your answer.
    \end{enumerate}
\end{problem}

\begin{solution}
    \bbni
    \begin{enumerate}
        \item [(a)] Note that $X, Y$ are independent events. Then the probability is, 

        \[\Pr(X  = Y) = \sum_{n=1}^\infty \Pr(X =n) \cdot \Pr(Y = n)\]
        
        Note that $X=n$ occurs if the first $n-1$ tosses are tails and the $n$-th toss is a head. Thus,
        \[ \Pr(X = n) =(1-p)^{\,n-1} p\]
        Similarly, 
        \[ \Pr(Y = n) =(1-p)^{\,n-1} p\]

        Therefore, 
        \begin{align*}
            \Pr(X  = Y) &= \sum_{n=1}^\infty \Pr(X =n) \cdot \Pr(Y = n) \\
            &= \sum_{n=1}^\infty \left( (1-p)^{\,n-1}  \right )^2 p^2 \\ 
            &= \sum_{n=1}^\infty  (1-p)^{\,2n-2}   p^2 \\ 
            &= p^2 \fr{1}{1 - (1-p)^2} \\ 
            &= \frac{p}{(2 - p)}
        \end{align*}

        Hence, the probability that $X=Y$ is $\dfrac{p}{2-p}$.

        \item [(b)] From the definition of expectation we know that:\\

        \begin{align*}
             E[Z] &= \sum_{k \geq 1} k \cdot Pr[Z=k] \\ 
             &= \sum_{k =1}^{\infty} k \cdot Pr[Z=k] \\ 
             &=  \sum_{k =1}^{\infty}\sum_{i=1}^{k}Pr[Z =k] \\ 
             &= \sum_{i=1}^\infty \sum_{k=i}^{\infty} Pr[Z=k]  \\ 
             &= \sum_{i=1}^\infty Pr[Z \geq i]
        \end{align*}

                    For each iteration of the inner sum we get, the sum of probabilities add up to  $P[Z \geq i]$. So we can write this following double sum as: 

                        \[
                        E[Z] = \sum_{i=1}^{\infty} Pr[Z \geq i] =  \sum_{i\geq1} Pr[Z \geq i] 
                         \]
                                


        
        \item [(c)] Let $Z = \min(X, Y)$. Then $E(Z)$
        \[E(Z) = \sum_{n=1}^\infty Pr[Z\geq n] \] from part b.
         We know:    
        \[Pr[Z \geq k]\] for an arbitrary variable denotes the event that the minimum number of tosses of X and Y is at least n. This means that both X and Y have had all tails on at least the first n-1 tosses. In other words: 
        \[Pr[Z \geq n] = Pr[X \geq n]\cdot Pr[Y \geq n]\]
        \[Pr[Z \geq n ] = (1-p)^{n-1} . (1-p)^{n-1}\]
        Plugging this in the above equation we get : 
        \begin{align*}
        E(Z) =& \sum_{n=1}^\infty Pr[Z\geq n] \\
             =& \sum_{n=1}^\infty(1-p)^{n-1} . (1-p)^{n-1} \\
             =& \sum_{n=1}^\infty(1-p)^{2n-2} 
        \end{align*}

        As this is a converging geometric series, the sum to infinity is give by:
        \begin{align*}
        E[Z] =& \frac{1}{1- (1-p)^{2}} \\
                    =& \frac{1}{(p)\cdot (2-p)}\\
        \end{align*}
    \end{enumerate} 
\end{solution}

\newpage

\begin{problem}{2}
    \textbf{Making friends.} 
    
    At the start of a new semester, $50$ new students at Uni High are assigned randomly to $50$ desks arranged in a $10 \times 5$ grid. Exactly half of the students come from middle school $A$ in the same city, and the other half coming from another middle school $B$ in one of the satellite cities.

    Every student has made up their mind: To talk to every student in the neighboring desks, but only if they came from a different middle school during the first day of class. In other words, a student from middle school $A$ who sits in the middle of the classroom will try to talk to each of the $4$ students at neighboring desks if they came from middle school $B$, and vice versa. Students who sit at the corner can only talk to at most two other students, and those who sit on the four sides may talk to three.
    
    Given that the student-to-desk assignment is completely random, how many pairs of students are expected to talk with each other?
\end{problem}

\begin{solution}
        \bbni 
        
        The total number of pairs is equal to 85. We can get this by counting the total numbers pairs along each of the 10 columns and 5 rows. Let X be a random variable denoting the total number of pairs of students are expected to talk with each other. 

        Let $X_k$ be an indicator variable indicating if the $k^{th}$ pair of students talk with each other. So we get, \[X = \sum_{k=1}^{85} X_k\]. 
        
         \[X_k = \begin{cases*}
            1 & \text{if students in the pair from different schools} \\ 
            0 & \text{otherwise}
        \end{cases*}\]

        
        As each $X_k$ is an indicator variable, $E[X_k]$  is given by the probability that each student in the $k^{th}$ pair is from a different middle school. This probability is given by

        \[Pr[X_k=1] = \frac{50 \cdot25}{50.49}  = \frac{25}{49}\]
        Hence, \[E[X_k] = \frac{25}{49}\]

        
        Finally, using linearity of expectations, we get: 
        \[E[X] = \sum_{k=1}^{85}E[X_k] = \frac{25}{49} \cdot85\]
\end{solution}

\newpage 

\begin{problem}{3}
    \textbf{Random permutations.} 
    
    Let $[n] = \{1,2,\dots,n\}$. Generate a permutation $\pi$ of $[n]$ uniformly at random. A fixed point of $\pi$ is an index $i$ such that $\pi(i)=i$.

    \begin{enumerate}
        \item[(a)] What is the expected number of fixed points? Justify your answer.

        Now start with the identity permutation on $[n]$. Then create a (not necessarily uniform) random permutation using the following algorithm: Let $t$ be a given nonnegative integer. Pick two distinct indices $i,j \in [n]$ uniformly at random and swap $\pi(i)$ with $\pi(j)$. Repeat this process $t$ times.
        
        Our goal is to mix the numbers well enough so that the resulting permutation, while not uniformly random, behaves like one at least in terms of the fixed point count.

        \item[(b)] Find the expected number of fixed points in $\pi$ after $t$ swaps. Justify your answer.

        \item[(c)] Let's say the permutation after $t$ swaps is quasi-random if the number of fixed points is $O(1)$.  How many swaps (asymptotically) do we need? That is, how large should t be?
    \end{enumerate}
\end{problem}


\begin{solution}
    \bbni
    \begin{enumerate}
        \item Let $X$ denote the expected number of fixed points, and $X_i$ for each $1 \leq i \leq n$ be defined as follows: 
        \[X_i = \begin{cases*}
            1 & \text{if $\pi(i) = i$} \\ 
            0 & \text{otherwise}
        \end{cases*}\]

        Then, $X = \sum_{i=1}^{n}X_i$, and by linearity of expectations, $E[X] = \sum_{i=1}^{n}E[X_i]$. Now note that, the probability that $\pi(i) = i$ is equal for each $i$, that is $E[X_i] = \frac{1}{n}$. 
        Then, 
        \[E[X] = \sum_{i=1}^{n} \frac{1}{n} = n \cdot \frac{1}{n} = 1\]
        That is, the expected number of fixed points is $1$. 

        \item 
        Let $p_t$ be the probability that $i$ is fixed after $t$ swaps. Since every index has uniform probability the expected number of fixed points after $t$ swaps $E[X_t] = np_t$. We can calculate $p_t$ as follows: 
        \[p_t = P[\text{fixed$|$ fixed before $t$-th swap}]p_{t-1} +  P[\text{fixed$|$ not fixed before $t$-th swap}](1-p_{t-1}) \]
        \begin{enumerate}
            \item If the $i$-th is fixed before the $t$-th swap: 
            
             Then $i$ would stay fixed if neither of the indices picked are $i$, that is the probability of it staying fixed is: 
             \[P[\text{fixed$|$ fixed before $t$-th swap}] = \frac{\binom{n-1}{2}}{\binom{n}{2}} = \frac{(n-1)!2! (n-2)!}{2! (n-3)! n!} = \frac{(n-2)}{n}\]

            \item If the $i$-th is not fixed before the $t$-th swap: 

            Then $i$ would be fixed if we the indices we pick are $i$ and the index that would fix $i$. The probability that $i$ gets picked is $2/n$, and the probability that we pick the index that fixes $i$ is $1/(n-1)$. The probability that the $t$-th swap fixes $i$ is: 
            \[P[\text{fixed$|$ not fixed before $t$-th swap}] = \frac{2}{n} \cdot \frac{1}{n-1} = \frac{2}{n(n-1)}\]
        \end{enumerate}

        That is the total probability of the $i$-th index being fixed after the $t$-th swap is: 
        \begin{align*}
            p_t &= P[\text{fixed$|$ fixed before $t$-th swap}]p_{t-1} +  P[\text{fixed$|$not fixed before $t$-th swap}](1-p_{t-1}) \\ 
            &=\frac{n-2}{n}p_{t-1} + \frac{2}{n(n-1)}(1- p_{t-1})  \\ 
            &= \frac{n-2}{n}p_{t-1} - \frac{2}{n(n-1)}p_{t-1} + \frac{2}{n(n-1)} \\ 
            &= \frac{(n-2)(n-1) - 2 }{n(n-1)} p_{t-1} + \frac{2}{n(n-1)}
        \end{align*}
        
        Since every index has uniform probability the expected number of fixed points after $t$ swaps $E[X_t] = np_t$, and after $t-1$ swaps $E[X_{t-1}] = np_{t-1}$. That gives us, 
        \[E[X_t] =  \frac{(n-2)(n-1) - 2 }{(n-1)} E[X_{t-1}] + \frac{2}{(n-1)} =  \frac{n-3}{n-1} E[X_{t-1}] + \frac{2}{n-1}\]

        For convenience let $\frac{n-3}{n-1} = a, \frac{2}{n-1} = b$. 
        
        \[E[X_{t}] =a E[X_{t-1}] +b\]

        
        Now note that by similar logic, 
        \[E[X_{t-1}] = aE[X_{t-2}] + b\]
        Substituting back we get, 
        \[E[X_{t}] =a (aE[X_{t-2}] + b) +b = a^2 E[X_{t-2}] + ab + b\]
        Where, 
        \[E[X_{t-2}] = a E[X_{t-3}] + b\]
        That is, 
        
        \[E[X_{t}] = a^2 (aE[X_{t-3}] + b) + ab + b = a^3 E[X_{t-3}] + a^2b + ab + b\]

        So on, until, 
        \[E[X_1] =a E[X_{0}] + b\]
        where $E[X_0] = n$. 
        Substituting this back in we get, 
        \[E[X_{t}] = a^t E[X_{0}] + a^{t-1}b + \cdots  + a^2b + ab + b =  a^t E[X_{0}] + b (a^{t-1} + \cdots  + a^2+ a + 1)\]
        Where the latter term is a geometric sum, 
        
        \[b (a^{t-1} + \cdots  + a^2+ a + 1) = b\fr{1 - a^{t}}{1 -a}\]

        \[\implies E[X_{t}] =  a^t E[X_{0}] + b \fr{1 - a^{t}}{ 1- a}\]

        Simplifying this, 
        \begin{align*}
            E[X_{t}] &=  a^t E[X_{0}] + b \fr{1 - a^{t}}{ 1- a} \\ 
            &= \fr{n-3}{n-1}^t n + \fr{2}{n-1} \fr{1 - \fr{n-3}{n-1}^t }{1 - \fr{n-3}{n-1}} \\ 
            &= \fr{n-3}{n-1}^t n + \fr{2}{n-1} \fr{1 - \fr{n-3}{n-1}^t }{\fr{2}{n-1}} \\ 
            &= \fr{n-3}{n-1}^t n + 1 - \fr{n-3}{n-2}^t  \\ 
            &= 1 + (n-1)\fr{n-3}{n-1}^t
        \end{align*}

        That is, the expected number of fixed points after $t$ swaps is 
        \[E[X_{t}] = 1 + (n-1)\fr{n-3}{n-1}^t\]

        \item For $E[X_t]$ to be $O(1)$ means that there exists some constant $c$ such that $T(n) \leq c$. That is, 
                \[(n-1)\fr{n-3}{n-1}^t \leq c\]
                

        \begin{align*}
            (n-1)\fr{n-3}{n-1}^t &\leq c \\
            \fr{n-3}{n-1}^t &\leq \frac{c}{n-1} \\
            t \log{\fr{n-3}{n-1}} &\leq \log c - \log {(n-1)}  \\ 
            t \log\left({1 - \frac{2}{n-1}}\right) &\leq \log c - \log {(n-1)}  \\ 
            t &\geq \frac{\log c - \log {(n-1)}}{ \log\left({1 - \frac{2}{n-1}}\right)}  \\ 
            t &\geq \frac{\log \fr{c} {n-1}}{ \log\left({1 - \frac{2}{n-1}}\right)}  
        \end{align*}
        Note that for sufficiently large $n$, $\log (1 - \frac{2}{n-1}) \approx - \frac{2}{n-1}$ so we get, 
        \[ t \geq \frac{n-1}{2}{\log \fr{n-1} {c}} \]
    \end{enumerate}
\end{solution}
\end{document}